% ============================================================
%  Chapter 1 — Introduction
% ============================================================
\section{Introduction}
\label{ch:introduction}

% ── Opening paragraphs — adapt from project plan introduction ──
The maritime shipping sector is undergoing a fundamental transition driven by
increasingly stringent environmental regulation. In recent years, regulatory
focus has shifted from long-term, fleet-level measures towards operational
performance, thereby directly influencing day-to-day decision-making. A
prominent example of this shift is the International Maritime Organization's
Carbon Intensity Indicator (\CII{}), which links a vessel's carbon dioxide
emissions to its transport work and assigns annual performance ratings based
on realized operational outcomes \parencite{imo2021cii}.

Unlike earlier regulatory instruments that primarily affected ship design or
fuel choice, the \CII{} framework directly constrains operational decisions
such as sailing speed, routing, and cargo acceptance. Consequently, shipping
operators are required to account for environmental performance already at the
planning stage, rather than evaluating emissions ex~post. This fundamentally
alters the nature of operational planning, as economic efficiency and
environmental compliance become tightly coupled \parencite{psaraftis2013speed}.

These challenges are particularly pronounced in short-sea shipping contexts
such as the Baltic Sea. The region is characterised by relatively short
sailing distances, frequent port calls, and flexible cargo patterns, which
increase the sensitivity of carbon intensity performance to operational
decisions. In such settings, traditional planning approaches that optimise
routing, speed, or cargo selection independently may lead to unintended
trade-offs and increased risk of regulatory non-compliance.

Adding to the \CII{} framework, the European Union extended its Emissions
Trading System to maritime transport from 2024, with full coverage of verified
emissions at EU ports applying from 2026 \parencite{eu_ets_shipping}.  The
\ETS{} attaches a direct financial cost to each tonne of CO$_2$ emitted, thereby
translating the environmental externality into a voyage cost item.  Together,
the \IMO{} \CII{} and the EU \ETS{} create a dual regulatory pressure on
Baltic Sea operators: the former through a reputational and compliance
mechanism, the latter through a direct cost mechanism.  Integrated operational
planning models that capture both instruments are therefore of immediate
practical relevance.

% ── Research gap ──────────────────────────────────────────────
\subsection{Research Gap}
\label{sec:intro:gap}

While recent academic literature has begun to address carbon-constrained
operational planning, there remains limited understanding of how carbon
intensity regulation affects integrated operational decisions in realistic
regional shipping networks. In particular, there is a lack of studies that
jointly consider routing, speed selection, and cargo acceptance under explicit
\CII{} constraints, and that translate such models into decision-support
insights relevant for operational planners \parencite{cheng2025cii}.

A further gap concerns the specifics of the \CII{} formula for Roll-on/Roll-off
(\RoRo{}) vessels. Unlike bulk or container shipping, where transport work is
measured in tonne-nautical-miles, the \IMO{} \CII{} formula for \RoRo{} vessels
uses Gross Tonnage (\GT{}) as the capacity measure rather than cargo weight.
This distinction has significant operational implications: cargo volume has no
direct effect on a \RoRo{} vessel's \CII{} rating, making speed and distance
the sole operational levers available to an operator.

% ── Research question ─────────────────────────────────────────
\subsection{Research Questions}
\label{sec:intro:rq}

This thesis addresses the identified gap by formulating and applying an
optimisation-based operational planning model for a short-sea \RoRo{} service
under \CII{} regulation. The central research question is:

\begin{quote}
\itshape
How does carbon intensity regulation reshape integrated routing and speed
decisions in short-sea shipping, and under what conditions does the
regulation become operationally binding?
\end{quote}

The following sub-questions guide the analysis:

\begin{enumerate}
  \item How should the \CII{} constraint be correctly formulated for a
    \RoRo{} vessel, and what are the operational implications of common
    misspecifications?
  \item Under what schedule and regulatory conditions does \CII{} compliance
    become infeasible, and what is the minimum structural change required
    to restore feasibility?
  \item What is the economic value of schedule flexibility for \CII{}
    compliance relative to a fixed timetable?
  \item How do post-2026 regulatory scenarios and \ETS{} carbon pricing
    interact to affect the cost structure of compliance?
\end{enumerate}

% ── Scope ─────────────────────────────────────────────────────
\subsection{Scope and Delimitations}
\label{sec:intro:scope}

The analysis is conducted within a stylised but operationally grounded
planning framework designed to capture structural decision
interdependencies rather than replicate company-internal planning systems.
The study is restricted to a single-vessel, short-sea shipping context,
with the Baltic Sea serving as the primary application environment.

The thesis does not address long-term strategic decisions such as fleet
renewal, vessel design, or fuel technology transitions. Carbon intensity
regulation is treated as an exogenous constraint under which operational
decisions are optimised, aligned with common modelling practice in recent
operational planning studies \parencite{cheng2025cii, psaraftis2013speed}.

% ── Contributions ─────────────────────────────────────────────
\subsection{Contributions}
\label{sec:intro:contributions}

The thesis makes the following contributions:

\begin{enumerate}
  \item A \MILP{} formulation for operational speed optimisation under
    \CII{} constraints, correctly parameterised for a \RoRo{} vessel
    including the \GT{}-based capacity measure and ice-class correction
    factors.
  \item A case study calibrated to \SCA{}'s M/V Obbola operating on a
    Baltic Sea loop, validated against EU \MRV{} operational data.
  \item A structured scenario analysis covering six regulatory and
    operational settings, including a demonstration of structural
    infeasibility under 2030 Rating~B targets.
  \item Quantification of the value of schedule flexibility: by replacing
    per-arc time budgets with a global cycle-time cap, the vessel reduces
    fuel consumption by 15.2\,\% (134.7~t per loop, EUR~2.12~million per year)
    while improving the attained \CII{} from 10.27 to 8.71~\gGTnm{}.
\end{enumerate}

% ── Thesis outline ────────────────────────────────────────────
\subsection{Thesis Outline}
\label{sec:intro:outline}

The remainder of this thesis is structured as follows.
\Cref{ch:literature} reviews the relevant academic and regulatory
literature. \Cref{ch:model} develops the mathematical optimisation model.
\Cref{ch:casestudy} describes the case study vessel and route.
\Cref{ch:scenarios} defines the scenario design. \Cref{ch:results}
presents and analyses the computational results. \Cref{ch:discussion}
discusses the findings, limitations, and managerial implications.
\Cref{ch:conclusion} concludes and identifies directions for future
research.
